% The article for the growth readout of the prediction ledger (GRW-1 to GRW-6).
\documentclass[11pt]{article}
\usepackage[T1]{fontenc}
\usepackage{lmodern}
\usepackage{microtype}
\usepackage[margin=1in]{geometry}
\usepackage{amsmath,amssymb,amsthm,mathtools}
\usepackage{algorithm}
\usepackage{algpseudocode}
\usepackage{booktabs}
\usepackage{enumitem}
\usepackage{url}
\usepackage[numbers,sort&compress]{natbib}
\usepackage[hidelinks]{hyperref}
\usepackage[capitalise,noabbrev]{cleveref}
% Distinct link targets for the line numbers of each algorithm.
\makeatletter
\providecommand{\theHALG@line}{}
\renewcommand{\theHALG@line}{\thealgorithm.\arabic{ALG@line}}
\makeatother

\theoremstyle{plain}
\newtheorem{theorem}{Theorem}[section]
\newtheorem{lemma}[theorem]{Lemma}
\newtheorem{proposition}[theorem]{Proposition}
\newtheorem{corollary}[theorem]{Corollary}
\theoremstyle{definition}
\newtheorem{definition}[theorem]{Definition}
\theoremstyle{remark}
\newtheorem{remark}[theorem]{Remark}

% A Lean declaration or a path, in monospace, breakable after an underscore, a dot or a slash.
\DeclareUrlCommand\lean{\def\UrlFont{\upshape\ttfamily}}
% The Lean declarations that check a theorem, set upright and ragged right.
\newenvironment{leannames}{\par\smallskip\small\raggedright\upshape\noindent Lean: }{\par}

\newcommand{\KL}[2]{\mathrm{KL}(#1\,\|\,#2)}
\newcommand{\E}{\mathbb{E}}
\newcommand{\R}{\mathbb{R}}
\newcommand{\abar}{\bar a}

\setlist{topsep=3pt,itemsep=1pt,parsep=0pt}
\renewcommand{\bibfont}{\small}
\emergencystretch=1.5em

\title{The Log Growth of a Proportional Stake Is the Cumulative Log-Score Advantage}
\author{Alberto Serrano-Calva}
\date{7 November 2026}

\begin{document}
\maketitle

\begin{abstract}
A pre-registered prediction ledger scores an arm's probability forecasts for yes/no questions
against a reference arm's forecasts on the same questions. We ask whether a sequence of stakes
informed by the forecasts can multiply a holding by a factor $N$, and what that requires. When the
reference's probability $r$ sets fair terms, the proportional stake multiplies the holding by $p/r$
on yes and by $(1-p)/(1-r)$ on no, for the arm's forecast $p$. We prove that its log growth over any
sequence is the arm's cumulative log-score advantage over the reference (GRW-1); that a question's
expected advantage under a true probability $q$ is $\KL{q}{r}-\KL{q}{p}$, largest for the truthful
forecast (GRW-2); that the stake maximizes expected log growth among constant fractions (GRW-3);
that the holding reaches $N$ times its start exactly when the cumulative advantage reaches $\log N$
(GRW-4); that the holding is never ruined (GRW-5); and that with independent, identically
distributed outcomes it tends to infinity almost surely when the expected advantage is positive,
and to zero when it is negative (GRW-6). All six are proved in Lean~4 with Mathlib, and the four
statements of the identity that the code relies on are judged against a separate specification.
The ledger's scorer reads its existing log scores as this growth; the readout is a computation over
scored pairs and stakes nothing.
\end{abstract}

\section{Introduction}

A forecast of a yes/no question can be scored once the question resolves. A pre-registered
prediction ledger records each forecast before its question closes and compares the forecasts of
one arm, a source of forecasts under a fixed configuration, with those of another on the
(question, round) pairs both predicted. It uses the Brier score~\citep{brier1950} and the
logarithmic score~\citep{good1952}, both strictly proper~\citep{gneiting2007}: the expected score
is best only at the true probability. A proper score measures how accurate the forecasts were.
This paper answers a different question about the same record: can a succession of stakes, each
informed by a forecast, multiply a starting holding by a large factor through many small
decisions rather than one large move, and what does such a path need?

The answer is a classical identity read against the ledger. Suppose that the reference arm's
probability $r$ for a question sets the terms of a stake on it: a unit placed on yes becomes $1/r$
units if the question resolves yes and nothing otherwise, and a unit placed on no becomes
$1/(1-r)$ units if it resolves no and nothing otherwise. The proportional
stake~\citep{kelly1956,cover2006} places the fraction $p$ of the holding on yes and $1-p$ on no,
for the arm's forecast $p$, and so multiplies the holding by the likelihood ratio of the outcome
under the arm's forecast and under the reference's. Over a sequence the factors multiply and their
logarithms add, so the logarithm of the growth is the sum, over questions, of the arm's log score
minus the reference's, both of which the ledger already computes. A path to $N$ times the start
exists exactly when this cumulative advantage reaches $\log N$.

The contributions are these.
\begin{itemize}
  \item A model, in Lean~4~\citep{moura2021} with Mathlib~\citep{mathlib2020}, of scored pairs,
    proportional and fractional stakes, the expectation over independent outcomes, and
    constant-fraction stakes, with six machine-checked theorems, GRW-1 to GRW-6
    (\cref{sec:theory}). The four statements of GRW-1, which the code relies on, are also judged
    by a comparator against a separate specification (\cref{sec:evaluation}).
  \item An account of why a score advantage is a growth rate, why the logarithmic score has this
    property and the Brier score does not, and what a fractional stake gives up and keeps: it is
    the proportional stake of a forecast moved toward the reference (\cref{sec:concepts}).
  \item A growth readout in the ledger's scorer, blinded like its primary comparison and read
    at the same $N_c$ clusters (\cref{sec:implementation}), with its tests, attacks and controls
    (\cref{sec:evaluation}).
\end{itemize}
The substance of GRW-1 to GRW-3 and of GRW-6 is classical~\citep{kelly1956,breiman1961,
thorp1969,cover2006}. What is new is the exact correspondence with a ledger's scored pairs, the
machine-checked statements and the readout. No result depends on observed data.

\section{Background and related work}
\label{sec:background}

\paragraph{The setting.} The ledger is an append-only, hash-chained record whose first entry, a
protocol written before any prediction, fixes the arms, the families of questions, the comparisons
and the number of clusters $N_c$ at which the primary comparison is read; each arm of the primary
comparison is pinned to the digest of its configuration before its first prediction, and questions
that turn on the same underlying event share a cluster. The scorer compares two arms by the paired
difference of Brier scores on the (question, round) pairs both predicted, with a bootstrap over
clusters or a sign-flip test over clusters, whichever the protocol declares, reports each arm's mean
Brier and log scores, checks controls, among them a recorded copy of the reference arm whose skill
must be exactly zero, and withholds the primary comparison until $N_c$ is fixed and reached and
every question of the first $N_c$ clusters is resolved or voided. From then on it reads those $N_c$
clusters only, the first whose pairs entered the ledger, however many the ledger holds.

\paragraph{Growth and information.} \citet{kelly1956} considered a gambler who learns of an outcome
through a noisy channel before betting on it at fair odds. The bet that maximizes the expected
logarithm of the holding places on each outcome the fraction of the holding equal to its
conditional probability, and the resulting exponential growth rate of the holding equals the rate
at which the channel transmits information. \citet{breiman1961} proved two asymptotic properties of
the strategy that maximizes the expected logarithm in repeated favorable games: the ratio of its
holding to that of any essentially different strategy tends to infinity almost surely, and it
asymptotically minimizes the expected time to reach a fixed large goal. \citet{thorp1969} surveyed
these results and their use in favorable games.

\paragraph{Gambling and relative entropy.} \citet[chap.~6]{cover2006} treat the horse race: with
outcome probabilities $p_i$, odds $o_i$ and fractions $b_i$ of the holding placed on the outcomes,
the doubling rate $W(b,p)=\sum_i p_i\log(b_i o_i)$ is maximized by proportional gambling, $b=p$.
With odds fair for a distribution $r$, $o_i=1/r_i$, it is $W(b,p)=D(p\,\|\,r)-D(p\,\|\,b)$, a
difference of relative entropies~\citep{kullback1951}, and side information $Y$ about the outcome
$X$ raises the optimal doubling rate by the mutual information $I(X;Y)$. GRW-2 is the two-outcome
case of the identity for fair odds, with the truth $q$ in place of $p$ and the arm's forecast in
place of $b$. Read per outcome rather than in expectation, the identity says that the logarithm of a
proportional stake's factor is a difference of two log scores, which is GRW-1.

\paragraph{Proper scoring rules.} \citet{brier1950} proposed scoring probability forecasts by the
mean squared difference between the forecast probabilities and the outcomes, and \citet{good1952}
proposed rewarding a probability by the logarithm of the probability it gave to what happened.
\citet{gneiting2007} review strictly proper scoring rules, characterize them through convex
functions, and associate with each a divergence, the expected shortfall of a forecast $p$ against
the truthful forecast when $q$ is true; for the logarithmic score it is $\KL{q}{p}$.

\paragraph{Formal developments.} \citet{carrone2026} proves, in Lean~4 with Mathlib, the Kelly
criterion for a bet that multiplies the stake by $1+u$ or by $1-d$ (its \lean{kelly_optimal} and
\lean{kelly_growth_pos}) and the almost-sure convergence of the time-average log growth (its
\lean{time_average_growth_ae}). The domain in the statement of GRW-3, the fractions that keep both
factors positive, and the shape of GRW-6, the strong law over pairwise independent, identically
distributed summands followed by the two limits, follow that development. The definitions and
proofs here are separate, over a model of reference, forecast and outcome in which the gain is
$u=1/r-1$ and the loss is $d=1$. Mathlib's strong law, \lean{strong_law_ae_real}, is Etemadi's
version~\citep{etemadi1981}, which needs only pairwise independence.

\section{Theory}
\label{sec:theory}

Logarithms are natural. The Lean names given with the definitions and theorems are in the
namespace \lean{Palace.Growth}, with the definitions in \lean{formal/Palace/Growth/Defs.lean} and
the theorems in \lean{formal/Palace/Growth.lean}, which contains no \lean{sorry}.

\subsection{Definitions}

\begin{definition}[Scored pair]
A \emph{scored pair} over a field $K$ is $x=(r,p)\in K^2$: the reference's probability $r$ that a
yes/no question resolves yes, and the arm's forecast $p$. An outcome is $o\in\{0,1\}$, with $o=1$
for yes, and the likelihood of a probability $y$ for $o$ is $L(y,1)=y$, $L(y,0)=1-y$. Over an
ordered field, $x$ is \emph{proper} if $0<r<1$ and $0<p<1$. (Lean: \lean{Pair}, \lean{likelihood},
\lean{Pair.Proper}.)
\end{definition}

The reference's probability sets the terms of a stake: a unit placed on yes becomes $1/r$ units on
yes and nothing on no, a unit placed on no becomes $1/(1-r)$ units on no and nothing on yes, and a
unit kept stays a unit; under the reference's probability every placement has expected multiplier
$1$, so the terms are fair.

\begin{definition}[Stakes, growth and holding]
\label{def:stake}
The \emph{proportional stake} places the fraction $p$ of the holding on yes and $1-p$ on no; its
factor is $F(x,o)=L(p,o)/L(r,o)$, that is $p/r$ on yes and $(1-p)/(1-r)$ on no. For a sequence
$\sigma=((x_1,o_1),\dots,(x_T,o_T))$ the \emph{growth} is $G(\sigma)=\prod_{i}F(x_i,o_i)$ and the
\emph{holding} from $W_0$ is $W(\sigma)=W_0\,G(\sigma)$. For $\lambda\in K$ the \emph{fractional
stake} has factor $F_\lambda(x,o)=1+\lambda(F(x,o)-1)$ and growth
$G_\lambda(\sigma)=\prod_iF_\lambda(x_i,o_i)$. (Lean: \lean{factor}, \lean{growth}, \lean{holding},
\lean{fractional}, \lean{fractionalGrowth}.)
\end{definition}

\begin{definition}[Scores, advantage, margin]
Over $\R$, the \emph{log score} of $y$ on $o$ is $s(y,o)=\log L(y,o)$ (the ledger's log loss is its
negative). The \emph{advantage} of a pair is $a(x,o)=s(p,o)-s(r,o)$, and the \emph{cumulative
advantage} of $\sigma$ is $A(\sigma)=\sum_ia(x_i,o_i)$. The \emph{margin} of a pair is
$m(x)=\min\{p,1-p,r,1-r\}$ and its \emph{worst factor} is $\min\{F(x,1),F(x,0)\}$. (Lean:
\lean{logScore}, \lean{advantage}, \lean{cumulative}, \lean{margin}, \lean{worstFactor}.)
\end{definition}

\begin{definition}[Truth, divergence, expectation]
A \emph{truth} for a pair is a probability $q$ that its outcome is yes. The \emph{expected
advantage} is $\abar_q(x)=q\,a(x,1)+(1-q)\,a(x,0)$ and the binary divergence is
$\KL{q}{y}=q\log(q/y)+(1-q)\log((1-q)/(1-y))$, with $0\log0=0$, which Lean's convention $\log0=0$
provides. A \emph{question} is a pair with a truth, $\xi=(x,q)$. For a list of questions with
independent outcomes, the expectation of a function $h$ of the outcomes is defined by recursion:
$\E_{[\,]}[h]=h([\,])$ and $\E_{\xi::\Xi}[h]=q\,\E_{\Xi}[h(1::\cdot)]+(1-q)\,\E_{\Xi}[h(0::\cdot)]$.
(Lean: \lean{expectedAdvantage}, \lean{klBin}, \lean{Question}, \lean{expect},
\lean{cumulativeOn}.)
\end{definition}

\begin{definition}[Constant fractions]
At reference $r$, the stake of a \emph{constant fraction} $f\in\R$ on yes places $f$ on yes and
keeps $1-f$: its factor is $Y(f)=1+f(1/r-1)$ on yes and $N(f)=1-f$ on no. Under a belief $p$ its
expected log growth is $g(f)=p\log Y(f)+(1-p)\log N(f)$, and the \emph{proportional fraction} is
$f^\star=(p-r)/(1-r)$. (Lean: \lean{stakeYes}, \lean{expectedLogGrowth},
\lean{proportionalFraction}.)
\end{definition}

A unit split $r:(1-r)$ between yes and no is worth one unit on either outcome, so for $p\ge r$
placing $p$ on yes and $1-p$ on no is the same stake as placing $f^\star$ on yes and keeping the
rest.

\subsection{The identity}

\begin{theorem}[GRW-1]
\label{thm:grw1}
\leavevmode
\begin{enumerate}
  \item[(i)] Over any field, $G(\sigma)=\prod_i L(p_i,o_i)\big/\prod_i L(r_i,o_i)$.
  \item[(ii)] Over any field, $G_1(\sigma)=G(\sigma)$.
  \item[(iii)] Over $\R$, if every pair of $\sigma$ is proper, $\log G(\sigma)=A(\sigma)$, and so
    $W(\sigma)=W_0\exp A(\sigma)$.
\end{enumerate}
\begin{leannames}
\lean{growth_eq_div}, \lean{fractionalGrowth_one}, \lean{log_growth}, \lean{holding_eq}, each
restated in \lean{formal/Spec/Growth.lean}.
\end{leannames}
\end{theorem}

\begin{proof}[Proof sketch]
(i) By induction: a product of quotients is the quotient of the products; Lean's division is total
($y/0=0$), and a zero likelihood of the reference makes both sides zero, so (i) needs no
hypothesis. (ii) unfolds the definitions. (iii) By induction: every factor of a proper pair is
positive (\cref{thm:grw5}), so the logarithm of the product is the sum of the logarithms of the
factors, and the logarithm of a factor is the advantage; then $\exp\circ\log$ applies.
\end{proof}

The identity is pathwise: it holds for every sequence of outcomes, with no probability model. The
readout computes the left side and reports the right side beside it: the difference of the two
arms' log scores, summed over the pairs both predicted.

\subsection{The expectation}

\begin{theorem}[GRW-2]
\label{thm:grw2}
Let $x=(r,p)$ be a proper pair.
\begin{enumerate}
  \item[(i)] For every real $q$, $\abar_q(x)=\KL{q}{r}-\KL{q}{p}$; so $\abar_q(x)>0$ if and only if
    $\KL{q}{p}<\KL{q}{r}$.
  \item[(ii)] (Gibbs) For $0\le q\le1$ and $0<y<1$, $\KL{q}{y}\ge0$, with $\KL{q}{y}>0$ when
    $y\ne q$; and $\KL{q}{q}=0$.
  \item[(iii)] For $0\le q\le1$, $\abar_q(x)\le\KL{q}{r}$, strictly when $p\ne q$, and
    $\abar_p(x)=\KL{p}{r}$: for a fixed truth, the truthful forecast has the largest expected
    advantage, and no other forecast attains it.
\end{enumerate}
\begin{leannames}
\lean{expectedAdvantage_eq}, \lean{expectedAdvantage_pos_iff}, \lean{klBin_nonneg},
\lean{klBin_pos}, \lean{klBin_self}, \lean{expectedAdvantage_le}, \lean{expectedAdvantage_lt},
\lean{expectedAdvantage_truth}.
\end{leannames}
\end{theorem}

\begin{proof}[Proof sketch]
(i) expands the logarithms of the quotients, with $q=0$ and $q=1$ treated apart. (ii) applies
$\log u\le u-1$ to $u=y/q$ and $u=(1-y)/(1-q)$ with weights $q$ and $1-q$; the weighted sum of
$u-1$ is zero, and the inequality is strict when $u\ne1$. (iii) combines (i) and (ii).
\end{proof}

By (i), $\abar_q(x)$ is the expected log score of the forecast minus the reference's: GRW-2 is the
strict propriety of the logarithmic score, read as growth. The edge it rewards is calibration, not
boldness.

\subsection{Optimality among constant fractions}

\begin{theorem}[GRW-3]
\label{thm:grw3}
Let $0<r<1$.
\begin{enumerate}
  \item[(i)] For $0<p<1$, $Y(f^\star)=p/r$ and $N(f^\star)=(1-p)/(1-r)$, both positive: at the
    proportional fraction the factors are the proportional stake's.
  \item[(ii)] For $0<p<1$ and every $f$ with $Y(f)>0$ and $N(f)>0$, $g(f)\le g(f^\star)$.
  \item[(iii)] $g(f^\star)=\KL{p}{r}$, which is positive when $0\le p\le1$ and $p\ne r$.
  \item[(iv)] For $0\le p\le1$, $g$ is concave on $(-r/(1-r),1)$, where both factors are positive;
    where neither vanishes, $g'(f)=p\,(1/r-1)/Y(f)-(1-p)/N(f)$, and for $0<p<1$,
    $g'(f^\star)=0$.
  \item[(v)] $f^\star\in[0,1)$ when $r\le p<1$, and $f^\star<0$ when $p<r$. A fraction $f$ on no, at
    the reference $1-r$ for no, has the factors of the fraction $-fr/(1-r)$ on yes and the same
    expected log growth; on no the proportional fraction is $(r-p)/r$, the maximizer among
    fractions on no that keep both factors positive, with the proportional stake's factors.
  \item[(vi)] $N(1)=0$: the whole holding on yes is lost on a no, which the logarithm does not
    represent.
\end{enumerate}
\begin{leannames}
(i) \lean{one_add_proportional}, \lean{one_sub_proportional}, \lean{stakeYes_proportional},
\lean{proportionalFraction_factors_pos}; (ii) \lean{expectedLogGrowth_le}; (iii)
\lean{expectedLogGrowth_proportional}, \lean{expectedLogGrowth_proportional_pos}; (iv)
\lean{concaveOn_expectedLogGrowth}, \lean{factors_pos_of_mem}, \lean{hasDerivAt_expectedLogGrowth},
\lean{hasDerivAt_expectedLogGrowth_proportional}; (v) \lean{proportionalFraction_mem},
\lean{proportionalFraction_neg}, \lean{stakeYes_mirror}, \lean{expectedLogGrowth_mirror},
\lean{proportionalFraction_mirror}, \lean{expectedLogGrowth_mirror_le},
\lean{stakeYes_mirror_proportional}; (vi) \lean{stakeYes_one_false}.
\end{leannames}
\end{theorem}

\begin{proof}[Proof sketch]
(i), (iii) and (v) are algebra, with (ii) applied at the reference $1-r$ and the belief $1-p$ for
the maximizer on no. (ii) does not use the derivative:
\begin{align*}
  g(f)-g(f^\star)&=p\log\frac{rY(f)}{p}+(1-p)\log\frac{(1-r)N(f)}{1-p}\\
  &\le p\Bigl(\frac{rY(f)}{p}-1\Bigr)+(1-p)\Bigl(\frac{(1-r)N(f)}{1-p}-1\Bigr)
  =rY(f)+(1-r)N(f)-1=0,
\end{align*}
by $\log u\le u-1$ and because $rY(f)+(1-r)N(f)=1$ for every $f$, which is the fairness of the
terms. (iv) composes the concave logarithm with affine maps and differentiates.
\end{proof}

\subsection{The path to N times the start}

\begin{theorem}[GRW-4]
\label{thm:grw4}
Let every pair of $\sigma$ be proper, and let $T$ be its length.
\begin{enumerate}
  \item[(i)] For $W_0>0$ and $N>0$, $W(\sigma)\ge N\,W_0$ if and only if $A(\sigma)\ge\log N$.
  \item[(ii)] Every factor of a pair lies in $[m(x),1/m(x)]$, and $|a(x,o)|\le\log(1/m(x))$.
  \item[(iii)] If $a(x_i,o_i)\le B$ for every $i$, then $A(\sigma)\le TB$; so if $B>0$ and
    $A(\sigma)\ge\log N$, then $T\ge\log N/B$.
  \item[(iv)] For a list $\Xi$ of $T$ questions with independent outcomes,
    $\E_\Xi[A]=\sum_i\abar_{q_i}(x_i)$. If this sum is $Ta$ with $a>0$, then $\E_\Xi[A]\ge\log N$
    if and only if $T\ge\log N/a$.
\end{enumerate}
\begin{leannames}
\lean{holding_ge_iff}, \lean{factor_mem}, \lean{abs_advantage_le}, \lean{cumulative_le},
\lean{length_ge_of_cumulative}, \lean{expect_cumulativeOn}, \lean{expect_const_add},
\lean{expect_ge_iff}.
\end{leannames}
\end{theorem}

\begin{proof}[Proof sketch]
(i) is \cref{thm:grw1}(iii) with the monotonicity of $\exp$. (ii) The likelihoods lie in
$[m(x),1]$, so their logarithms lie in $[\log m(x),0]$. (iii) sums the bound. (iv) By induction on
the list: the two weights of a question sum to one, so a constant passes through the expectation,
and the expectation of a sum is the sum of the expectations.
\end{proof}

At a mean expected advantage of $a=0.05$ nats per question, (iv) gives $\log10/0.05\approx46.05$,
so 47 questions for the expected cumulative advantage to reach $\log10$; this concerns the
expectation of the logarithm, not the probability of having reached ten times the start
(\cref{sec:dispersion}). Part (iii) is a floor that holds whatever the outcomes.

\subsection{No ruin}

\begin{theorem}[GRW-5]
\label{thm:grw5}
Over an ordered field, let every pair of $\sigma$ be proper.
\begin{enumerate}
  \item[(i)] Every factor is positive; so $G(\sigma)>0$ and, for $W_0>0$, $W(\sigma)>0$.
  \item[(ii)] The worst factor of $x$ is $\min\{p/r,(1-p)/(1-r)\}$; it is positive, at most either
    factor, and at least $m(x)$.
  \item[(iii)] If $\ell\ge0$ is at most the worst factor of every pair of $\sigma$, then
    $\ell^T\le G(\sigma)$.
\end{enumerate}
\begin{leannames}
\lean{factor_pos}, \lean{growth_pos}, \lean{holding_pos}, \lean{worstFactor_eq},
\lean{worstFactor_pos}, \lean{worstFactor_le}, \lean{margin_le_worstFactor}, \lean{pow_le_growth}.
\end{leannames}
\end{theorem}

\begin{proof}[Proof sketch]
Quotients of positive numbers are positive; (ii) uses \cref{thm:grw4}(ii); (iii) is an induction in
which $\ell$ lies below each realized factor.
\end{proof}

The fractional stake keeps (i) for $0<\lambda\le1$, where $F_\lambda=(1-\lambda)+\lambda F$ is a
positive combination, and loses it beyond: at $\lambda=2$ a factor below $1/2$ becomes negative.

\subsection{The time average}

\begin{theorem}[GRW-6]
\label{thm:grw6}
Let $(\Omega,P)$ be a probability space and $o_0,o_1,\ldots\colon\Omega\to\{0,1\}$ pairwise
independent outcomes, each yes with probability $q$. For a pair $x$, let $A_n(\omega)$ be the
cumulative advantage of $x$ repeated $n$ times with the outcomes
$o_0(\omega),\dots,o_{n-1}(\omega)$, and $W_n(\omega)$ the holding after them, from $W_0$.
\begin{enumerate}
  \item[(i)] Almost surely, $A_n/n\to\abar_q(x)$.
  \item[(ii)] If $x$ is proper, $W_0>0$ and $\abar_q(x)>0$, then almost surely $W_n\to\infty$, so
    the holding eventually exceeds $N\,W_0$, for every $N$.
  \item[(iii)] If $x$ is proper and $\abar_q(x)<0$, then almost surely $W_n\to0$.
\end{enumerate}
\begin{leannames}
\lean{time_average_advantage_ae}, \lean{holding_tendsto_atTop_ae},
\lean{holding_tendsto_zero_ae}, with Mathlib's \texttt{bernoulliMeasure true false q} as the law of
each outcome.
\end{leannames}
\end{theorem}

\begin{proof}[Proof sketch]
The advantages $X_i=a(x,o_i)$ take two values, so they are integrable, and as one measurable
function of the $o_i$ they are pairwise independent and identically distributed. Mathlib's
\lean{strong_law_ae_real} gives $A_n/n\to\E[X_0]=\abar_q(x)$ almost surely. A sequence whose
average tends to a positive limit tends to $+\infty$, one whose average tends to a negative limit
tends to $-\infty$, and $W_n=W_0\exp A_n$ by \cref{thm:grw1}(iii).
\end{proof}

\section{Concepts and intuition}
\label{sec:concepts}

The two propositions of this section are elementary consequences of the theorems; they are not
among the machine-checked statements.

\subsection{A score advantage is a growth rate}

Under fair terms, every stake is a forecast. Split the holding into a part $y$ placed on yes, $z$
on no and $k$ kept, nonnegative with $y+z+k=1$. The factors are $\alpha=y/r+k$ on yes and
$\beta=z/(1-r)+k$ on no, and $r\alpha+(1-r)\beta=1$. With $\hat p=r\alpha\in[0,1]$,
$\alpha=\hat p/r$ and $\beta=(1-\hat p)/(1-r)$: the factors of the proportional stake of the
forecast $\hat p$. Every stake that places no more than the holding is therefore the proportional
stake of the forecast it implies, and by \cref{thm:grw1} its log growth is that forecast's
log-score advantage over the reference. To choose a stake is to choose a forecast. Factors
multiply and logarithms add, so no single question has to contribute much: a path to $N$ times the
start is a succession of small edges, each a question on which the arm gave the outcome more
probability than the reference did. An arm that copies the reference places $r$ on yes and $1-r$ on
no, which is the same as keeping everything, and grows by exactly $1$.

\subsection{Why the logarithmic score and not the Brier score}

Both scores are strictly proper, and the ledger's pre-registered comparison uses the Brier score.
Only the logarithmic score's differences are log growth. The log factor of a stake that implies
$\hat p$ is $\log(\hat p/r)$ on yes and $\log((1-\hat p)/(1-r))$ on no. A score $S$ with
$S(\hat p,o)-S(r,o)$ equal to this for all $\hat p,r\in(0,1)$ has $S(\hat p,1)-\log\hat p$
independent of $\hat p$, and likewise $S(\hat p,0)-\log(1-\hat p)$: it is the logarithmic score up
to a constant for each outcome. The Brier score is bounded, so its differences cannot follow a
factor that a confident error makes arbitrarily small.

A constructed example shows the two disagreeing in sign. Ten questions all resolve yes; the
reference gives each $1/2$; the arm forecasts $0.9$ on nine and $0.001$ on the tenth. The arm's
mean Brier score is $0.109$ against the reference's $0.25$. Its cumulative log-score advantage is
$9\log1.8+\log0.002\approx5.290-6.215=-0.925$, so the proportional stake ends at
$e^{-0.925}\approx0.397$ of its start. On the tenth question the Brier scores differ by less than
$1$, while the factor is $1/500$. A positive Brier advantage does not imply growth; the readout and
the primary comparison answer different questions, and the scorer reports both.

\subsection{Fractional stakes}

\begin{proposition}[A fractional stake is a forecast moved toward the reference]
\label{prop:shrink}
For a pair $x=(r,p)$ with $r\notin\{0,1\}$ and any $\lambda$, $F_\lambda(x,o)=F(x_\lambda,o)$, where
$x_\lambda=(r,\,r+\lambda(p-r))$.
\end{proposition}

\begin{proof}
On yes, $(r+\lambda(p-r))/r=1+\lambda(p/r-1)$; on no,
$(1-r-\lambda(p-r))/(1-r)=1+\lambda\bigl((1-p)/(1-r)-1\bigr)$.
\end{proof}

For $0\le\lambda\le1$ and a proper pair, $r+\lambda(p-r)$ lies between $r$ and $p$, so $x_\lambda$
is proper and \cref{thm:grw1} applies to it: the $\lambda$-stake's log growth is the cumulative
log-score advantage of the forecast $r+\lambda(p-r)$. This is the sense in which the identity holds
at $\lambda=1$ only: at other stakes the growth is the advantage of a different forecast, not the
arm's. By \cref{thm:grw2}(i) applied to $x_\lambda$, under a truth $q$,
\begin{equation}
  \label{eq:fractional}
  \E_q[\log F_\lambda]=\KL{q}{r}-\KL{q}{r+\lambda(p-r)}.
\end{equation}

\begin{proposition}
\label{prop:fraction}
Let $x=(r,p)$ be proper, $0\le q\le1$ and $0\le\lambda\le1$.
\begin{enumerate}
  \item[(i)] $\E_q[\log F_\lambda]\ge\lambda\,\E_q[\log F_1]$.
  \item[(ii)] At $\lambda=0$, $\frac{d}{d\lambda}\E_q[\log F_\lambda]=\E_q[F-1]
    =(q-r)(p-r)/(r(1-r))$.
  \item[(iii)] If $0<q<1$ and $\lambda^\star=(q-r)/(p-r)\in(0,1]$, then $\lambda^\star$ maximizes
    $\E_q[\log F_\lambda]$ over $[0,1]$, with maximum $\KL{q}{r}$.
\end{enumerate}
\end{proposition}

\begin{proof}
(i) $y\mapsto\KL{q}{y}$ is convex, so
$\KL{q}{r+\lambda(p-r)}\le(1-\lambda)\KL{q}{r}+\lambda\KL{q}{p}$; substitute in
\eqref{eq:fractional} and use \cref{thm:grw2}(i). (ii) differentiates $\log(1+\lambda(F-1))$ at $0$
and expands $q\,p/r+(1-q)(1-p)/(1-r)-1$. (iii) By
\cref{thm:grw2}(iii) the expected advantage is largest when the forecast is the truth,
$r+\lambda(p-r)=q$.
\end{proof}

A fractional stake keeps at least the fraction $\lambda$ of the full stake's expected log growth,
whatever the truth; for small $\lambda$ it grows in expectation whenever the forecast and the truth
lie on the same side of the reference, a weaker condition than $\KL{q}{p}<\KL{q}{r}$; and an arm
whose forecasts lie twice as far from the reference as the truth does best at the half stake
(\cref{tab:fraction}). Near the reference, with $\chi^2=(p-r)^2/(r(1-r))$, a truthful forecast has
$\E_q[\log F_\lambda]\approx(\lambda-\lambda^2/2)\chi^2$ and a standard deviation of
$\log F_\lambda$ near $\lambda\chi$: the half stake keeps about three quarters of the growth rate
with half the spread.

\begin{table}[t]
\centering
\caption{Expected log growth per question, from \eqref{eq:fractional}, and the standard deviation
of the log factor, for a reference $r=0.5$ and a truth $q=0.55$. The values describe the formulas,
not any data.}
\label{tab:fraction}
\small
\begin{tabular}{cccc}
\toprule
forecast $p$ & stake $\lambda$ & $\E_q[\log F_\lambda]$ (nats) & standard deviation \\
\midrule
$0.55$ & $1$   & $0.00501$  & $0.0998$ \\
$0.55$ & $1/2$ & $0.00375$  & $0.0498$ \\
$0.55$ & $1/4$ & $0.00219$  & $0.0249$ \\
$0.60$ & $1$   & $-0.00014$ & $0.2017$ \\
$0.60$ & $1/2$ & $0.00501$  & $0.0998$ \\
$0.60$ & $1/4$ & $0.00375$  & $0.0498$ \\
\bottomrule
\end{tabular}
\end{table}

\subsection{Expectation and dispersion}
\label{sec:dispersion}

With independent questions of mean $\mu$ and standard deviation $s$ per question, the cumulative
advantage after $T$ questions has mean $T\mu$ and standard deviation $s\sqrt T$, so the drift
exceeds one standard deviation only once $T>(s/\mu)^2$. For the truthful full stake of
\cref{tab:fraction}, $(s/\mu)^2\approx397$ and $\log10/\mu\approx460$: under a normal approximation
the holding is above ten times its start about half the time at question 460, the count
\cref{thm:grw4}(iv) gives. The theorems say when a path exists and how long it is for a given edge;
they do not say that an arm has an edge, which is what the pre-registered comparison measures.

\section{Implementation}
\label{sec:implementation}

The readout is in \lean{eval/predictions/scoring.py} (\lean{factor}, \lean{growth_path},
\lean{growth}, \lean{growth_report}) and runs as
\texttt{python -m eval.predictions growth -{}-ledger <file> [-{}-json] [-{}-unblind]}; its design is
section G8 of \lean{docs/design-notes/dn-prediction-ledger.md}. From the protocol it takes the
reference arm, the primary comparison, the clip $c$ of the log score (default $0.001$), the target
$N$ (\lean{growth_target}, default $10$, which must exceed $1$) and $N_c$.

Both probabilities are clipped to $[c,1-c]$, as the log score clips them, so every pair is proper,
the identity is with the ledger's own clipped log score, and every factor lies in
$[c/(1-c),(1-c)/c]$. For an arm and a family the pairs are those both the arm and the reference
predicted, in the order their outcomes became known, ties broken by question and round; the growth
and the geometric mean do not depend on the order, the drawdown does. For each arm other than the
reference, each family and each stake $\lambda\in\{1,1/2,1/4\}$, the readout reports the number of
pairs, the growth, the cumulative advantage, the per-pair geometric mean, the worst factor, the
drawdown (the largest fall of the running product from its running peak, as a fraction of the
peak), and the further pairs needed to reach $N$ at the realized mean log growth per pair. The
stake must lie in $(0,1]$; beyond $1$ a factor can be negative.

Until $N_c$ is fixed, the primary comparison has that many clusters and every question of the first
$N_c$ is resolved or voided, its two arms are withheld on its family and named with the reason,
unless the caller passes \lean{--unblind}. Other arms are shown, among them the recorded copy of
the reference, whose factors are all exactly $1$. Once $N_c$ is fixed the readout reads only the
first $N_c$ clusters of the primary comparison, in the order their first pair entered the ledger,
which is the cut the report makes, and says how many clusters and pairs it left out; a reader who
waits past $N_c$ reads what one who stopped there reads. Once the comparisons are readable, the full
report includes the readout. \Cref{alg:path,alg:report} give the
procedure; it costs a sort per arm and family and a linear pass per stake.

\begin{algorithm}[t]
\caption{The growth path of one arm at one stake.}
\label{alg:path}
\small
\begin{algorithmic}[1]
\Require factors $F_1,\dots,F_n$ with $n\ge1$, a stake $\lambda\in(0,1]$, a target $N>1$
\State $W\gets1$; $P\gets1$; $D\gets0$
\For{$i=1,\dots,n$}
  \State $S_i\gets1+\lambda(F_i-1)$; \quad $W\gets W\cdot S_i$; \quad $P\gets\max(P,W)$;
    \quad $D\gets\max(D,1-W/P)$
\EndFor
\State $\hat\mu\gets(\log W)/n$
\If{$W\ge N$} $\tau\gets0$ \ElsIf{$\hat\mu>0$} $\tau\gets\lceil(\log N-\log W)/\hat\mu\rceil$
  \Else{} $\tau\gets$ none \EndIf
\State \Return growth $W$, geometric mean $e^{\hat\mu}$, worst factor $\min_iS_i$, drawdown $D$,
  further pairs $\tau$
\end{algorithmic}
\end{algorithm}

\begin{algorithm}[t]
\caption{The growth readout of a ledger.}
\label{alg:report}
\small
\begin{algorithmic}[1]
\Require a scored table; a protocol with reference $R$, primary comparison (family $\phi_0$, arm
  $\alpha$, baseline $\beta$), clip $c$, target $N$; a flag \textit{unblind}
\State $\textit{reason}\gets$ none if \textit{unblind}, or if $N_c$ is fixed and reached and every
  question of the first $N_c$ clusters is settled; else why not
\If{$N_c$ is fixed} keep the first $N_c$ clusters of the primary comparison, by entry \EndIf
\For{each family $\phi$ and each arm $A\ne R$}
  \If{$\textit{reason}\ne$ none \textbf{and} $\phi=\phi_0$ \textbf{and} $A\in\{\alpha,\beta\}$}
    record $A$ on $\phi$ as withheld, and skip it
  \EndIf
  \State $K\gets$ the pairs $A$ and $R$ both predicted on $\phi$, by time of resolution, then key
  \State for $k\in K$: $p_k,r_k\gets$ the two forecasts clipped to $[c,1-c]$, $F_k\gets
    L(p_k,o_k)/L(r_k,o_k)$
  \For{$\lambda\in(1,1/2,1/4)$}
    emit $A$, $\phi$, $\lambda$, $|K|$, $\sum_k\bigl(s(p_k,o_k)-s(r_k,o_k)\bigr)$ and
    \Call{GrowthPath}{$F,\lambda,N$}
  \EndFor
\EndFor
\State \Return the readouts, the withheld arms, \textit{reason} and what the cut left out
\end{algorithmic}
\end{algorithm}

\section{Evaluation}
\label{sec:evaluation}

\subsection{What is machine-checked}

The development uses Lean~4 at v4.36.0-rc2 with Mathlib at the same tag: 129 lines of definitions,
702 lines with 67 theorems, a specification of 59 lines and 88 lines of attacks. The four
statements of GRW-1 are restated in \lean{formal/Spec/Growth.lean} with their proofs left as
\lean{sorry}, and the comparator checks that the library proves each exactly as stated, from the
axioms \lean{propext}, \lean{Classical.choice} and \lean{Quot.sound} alone; with this change it
judges 285 statements across the library. GRW-2 to GRW-6 are design theorems, proved in the same
library and checked by the kernel but not restated in the specification, because no code relies on
them. The gate, \lean{formal/check.sh}, builds with warnings as errors, runs the attacks, holds the
count of \lean{sorry} at zero, audits every declaration's axioms, replays the exported library
through Lean's kernel and five further checkers, and runs the comparator. Run on the commit
that introduced the readout, it reported 4353 declarations audited, none using a non-standard
axiom or depending on \lean{sorry}, 4308 exported and replayed, and 285 statements
judged.

\subsection{Attacks and controls}

The laws are attacked with Plausible~\citep{plausible} over $\mathbb{Q}$, where the logarithms do
not evaluate, so GRW-1 is attacked in its product form. A pair is built from four
natural numbers as $r=(a+1)/(a+b+2)$ and $p=(c+1)/(c+d+2)$, so every pair is proper, and a sequence
is a list of such pairs with outcomes. An attack fails the build on a counterexample or when fewer
than 100 instances pass its guards; a control states a natural mistake and fails the build unless a
counterexample is found. The seeds are fixed, so \cref{tab:attacks} is reproducible: it is printed
by \texttt{lake env lean PalaceTest/Growth.lean} in \lean{formal/} and appears unchanged in the
continuous-integration log. The control at seed 329, for one, fails at $r=1/2$, $p=1/5$ and a yes,
where the factor $2/5$ becomes $-1/5$ at $\lambda=2$.

\begin{table}[t]
\centering
\caption{The attacks, of 500 instances each, and the controls on the growth laws.}
\label{tab:attacks}
\small
\begin{tabular}{rll}
\toprule
seed & statement & outcome \\
\midrule
320 & GRW-1: the growth is the likelihood ratio & survived: 500 tested, 0 discarded \\
321 & GRW-1 without the division & survived: 500 tested, 0 discarded \\
322 & control: the growth is the arm's likelihood alone & refuted after 5 instances \\
323 & control: the half stake satisfies the identity & refuted after 6 instances \\
324 & GRW-3: at $f^\star$ the factors are the proportional stake's & survived: 500 tested,
  0 discarded \\
325 & control: the fraction $p$ is the proportional stake & refuted on the first instance \\
326 & GRW-4: every factor lies in $[m,1/m]$ & survived: 500 tested, 0 discarded \\
327 & control: every factor is at most $1/r$ & refuted after 30 instances \\
328 & GRW-5: $G>0$ and $\ell^T\le G$ & survived: 500 tested, 0 discarded \\
329 & control: the stake $\lambda=2$ keeps the growth positive & refuted after 18 instances \\
\bottomrule
\end{tabular}
\end{table}

\subsection{The identity on random ledgers}

The code that relies on GRW-1 is tested in \lean{tests/unit/test_prediction_scoring.py}. The
identity test draws, with Hypothesis~\citep{maciver2019}, 40 ledgers of 1 to 25 rows (a reference
probability and a forecast from $[0,1]$, endpoints included so that the clip is exercised, and an
outcome), writes each through the ledger's own interface and reads it back through the scorer: each
factor is the exponential of its pair's advantage and the growth the exponential of the cumulative
advantage, within a relative $10^{-9}$, and the growth and the worst factor are positive. Its
control mutates the factor: the half stake's growth and a product read without the reference both
differ from the exponential of the cumulative advantage. Further tests check a hand sequence
(factors $2$, $1/2$, $2$: growth $2$, drawdown $1/2$, 7 further pairs to ten times the start), that
the copy of the reference grows by exactly $1$ at every stake, the blinding, the cut at $N_c$ and
the command line. The file's 46 tests pass (\texttt{uv run pytest -q
tests/unit/test\_prediction\_scoring.py}), as does the repository's full selection on this branch
(3370 passed, 11 skipped). None of this is a claim about any arm's forecasts.

\section{Limitations}
\label{sec:limitations}

\begin{itemize}
  \item The readout is a computation over scored pairs. It places nothing, holds nothing and acts on
    nothing; its stake is a thought experiment that reads the ledger's log scores as a holding.
  \item The model assumes fair terms at the reference's probability on every question, for any
    amount, with no costs and no spread between the terms for yes and for no; a cost of a fraction
    $\kappa$ of the holding per question subtracts $-\log(1-\kappa)\approx\kappa$ from every log
    factor and can turn a positive advantage negative. The reference does not respond to the stake.
  \item GRW-6 assumes one pair repeated with pairwise independent, identically distributed outcomes,
    and GRW-4(iv) independent outcomes; the ledger's questions differ, and those that share a
    cluster are dependent by construction. GRW-1 is pathwise and holds regardless.
  \item The truth $q$ is a modelling device that nothing estimates, and no finite-time probability
    of reaching $N$ times the start is proved.
  \item Only GRW-1 is judged against a separate specification. The binary divergence is written out
    rather than derived from Mathlib's divergence of measures, so its agreement with the divergence
    of two Bernoulli laws rests on inspection of the definition, not on a proved bridge.
  \item The readout multiplies in double precision: with the default clip, 103 consecutive factors
    at the top bound ($999$) overflow the product, and 108 at the bottom bound underflow it to zero,
    where its logarithm raises an error. Both need the forecast and the reference at opposite clips
    on consecutive questions; summing logarithms, the same quantity by GRW-1, avoids both.
\end{itemize}

\section{Future work}
\label{sec:future}

\paragraph{The strong law without identical distribution.} With margins bounded below by $m$,
every advantage is bounded by $\log(1/m)$ (\cref{thm:grw4}(ii)), so for independent questions that
differ Kolmogorov's strong law gives $(A_n-\E A_n)/n\to0$ almost surely. Mathlib's strong law at
this version is for identically distributed summands; a formal version for bounded independent
summands would extend GRW-6 to the ledger's varied questions.

\paragraph{Growth under costs.} When the terms for yes and for no imply probabilities that sum to
more than one, the stake that maximizes expected log growth on one question places nothing when the
forecast lies between them; extending GRW-1 and GRW-3 to such terms would let the readout report
growth under a declared transaction cost.

\paragraph{Side information as mutual information.} If the reference is calibrated to its
information and the arm to information that contains the reference's, the arm's expected advantage
is the conditional mutual information between the outcome and the arm's information given the
reference's: the two-outcome case of the side-information result of \citet{kelly1956} and
\citet[chap.~6]{cover2006}, and a statement of what the readout measures.

\paragraph{Smaller steps.} A proved bridge from the written-out divergence to Mathlib's; questions
with more than two outcomes; a finite-time bound by Hoeffding's inequality for advantages bounded by
$\log(1/m)$; and a readout that accumulates logarithms.

\section{Conclusion}

Against a reference whose probabilities set fair terms, the log growth of the proportional stake is
the cumulative log-score advantage, exactly and for every sequence of outcomes; its expectation is a
difference of divergences, it is the best constant fraction, it is never ruined, and with
independent repetitions the sign of the expected advantage decides between growth without bound and
decay to zero. These statements are machine-checked, and the ledger's scorer reads its log scores
as this holding. Whether an arm has the edge such a path needs is the question the pre-registered
comparison answers.

\bibliographystyle{plainnat}
\bibliography{refs}
\end{document}
